% Session 7 -- MERA and PEPS: the two architectures the section names but never draws.
% Placement: sec. "Beyond MPS: MERA and PEPS", after the PEPS paragraph.
%
% LAYOUT NOTE: pure TikZ, canvas coordinates; one-off tensor shapes declared with
% \tikzset inside this picture only.
\begin{aqmfigure}{fig:s07:merapeps}{The two tensor networks beyond the chain, and what their geometry
buys. Left: the multi-scale entanglement renormalization ansatz. Its tensors are stacked in layers,
each layer first removing the short-range entanglement between neighbouring blocks (the square
disentanglers) and then coarse-graining two sites into one (the triangular isometries), so that the
network describes the state at successively longer length scales rather than at successive positions
along a chain. A bipartition now cuts one bond in every layer instead of one bond in total, which is
precisely how MERA supports the logarithmically growing entropy (7.12) of a critical chain that an
MPS of fixed bond dimension cannot reach. Right: projected entangled pair states, the honest 2D
generalization of (7.5), with one tensor per lattice site carrying a physical leg and four virtual
bonds. The picture also shows the difficulty: there is no way to sweep along such a network one
tensor at a time with the rest reduced to a simple environment, and contracting it exactly is itself
a hard problem, which is why PEPS has no clean analogue of DMRG.}

\begin{tikzpicture}[x=1cm, y=1cm]
\tikzset{
  tens/.style={aqm, circle, draw=black, fill=aqmfill, minimum size=6.5pt, inner sep=0pt},
  diseng/.style={aqm, rectangle, draw=black, fill=white, minimum size=8pt, inner sep=0pt},
  isom/.style={aqm, regular polygon, regular polygon sides=3, draw=black, fill=aqmfill,
               minimum size=11pt, inner sep=0pt},
  leg/.style={aqm, draw=black, line width=\aqmthinwidth}
}

% ============================ (a) MERA ======================================
\node[aqmpanel, anchor=west] at (-0.10,3.30) {(a) MERA: layers across length scales};

% --- physical sites -----------------------------------------------------------
\foreach \i in {0,...,7} {
  \node[tens] (p\i) at (0.42*\i,0) {};
  \draw[leg] (0.42*\i,-0.08) -- (0.42*\i,-0.34);
}

% --- layer 1: disentanglers on (1,2), (3,4), (5,6) ---------------------------
\foreach \c/\l/\r in {0.63/1/2, 1.47/3/4, 2.31/5/6} {
  \node[diseng] (u\l) at (\c,0.62) {};
  \draw[leg] (p\l) -- (u\l);
  \draw[leg] (p\r) -- (u\l);
}
% --- layer 1: isometries on (0,1), (2,3), (4,5), (6,7) -----------------------
\foreach \c/\n in {0.21/0, 1.05/1, 1.89/2, 2.73/3} {
  \node[isom] (w\n) at (\c,1.32) {};
}
\draw[leg] (p0) -- (w0);
\draw[leg] (u1) -- (w0);
\draw[leg] (u1) -- (w1);
\draw[leg] (u3) -- (w1);
\draw[leg] (u3) -- (w2);
\draw[leg] (u5) -- (w2);
\draw[leg] (u5) -- (w3);
\draw[leg] (p7) -- (w3);

% --- layer 2 -----------------------------------------------------------------
\node[diseng] (uu) at (1.47,1.92) {};
\draw[leg] (w1) -- (uu);
\draw[leg] (w2) -- (uu);
\node[isom] (v0) at (0.63,2.56) {};
\node[isom] (v1) at (2.31,2.56) {};
\draw[leg] (w0) -- (v0);
\draw[leg] (uu) -- (v0);
\draw[leg] (uu) -- (v1);
\draw[leg] (w3) -- (v1);
\node[tens] (top) at (1.47,3.02) {};
\draw[leg] (v0) -- (top);
\draw[leg] (v1) -- (top);

% --- the scale axis ----------------------------------------------------------
\draw[aqmarrowmuted] (-0.62,-0.20) -- (-0.62,3.00);
\node[aqmlabmuted, rotate=90, anchor=south] at (-0.76,1.40) {coarser length scale};

% --- legend and message ------------------------------------------------------
\node[diseng] at (3.75,0.62) {};
\node[aqmlabsmall, anchor=west] at (3.92,0.62) {disentangler};
\node[isom] at (3.75,1.32) {};
\node[aqmlabsmall, anchor=west] at (3.92,1.32) {isometry ($2\to1$ sites)};
\node[aqmlabsmall, anchor=west, align=left, text width=3.3cm, text=aqmaccent] at (3.65,2.45)
  {a bipartition cuts one bond \emph{per layer}, so $S$ can grow like $\log\ell$: the critical
   law (7.12)};
\node[aqmlabsmall, anchor=north, text=aqmmuted] at (1.47,-0.48) {physical legs};

% ============================ (b) PEPS ======================================
\node[aqmpanel, anchor=west] at (7.30,3.30) {(b) PEPS: a 2D grid of tensors};

\begin{scope}[shift={(8.20,0.85)}]
  \foreach \i in {0,1,2,3} { \foreach \j in {0,1,2,3} {
      \node[tens] (t\i\j) at (0.62*\i,0.62*\j) {};
      \draw[leg] (0.62*\i-0.05,0.62*\j-0.05) -- (0.62*\i-0.24,0.62*\j-0.24);
  }}
  \foreach \i/\k in {0/1, 1/2, 2/3} { \foreach \j in {0,1,2,3} {
      \draw[leg] (t\i\j) -- (t\k\j);
      \draw[leg] (t\j\i) -- (t\j\k);
  }}
\end{scope}
\node[aqmlabsmall, anchor=north, align=center, text width=4.4cm, text=aqmmuted] at (9.13,0.44)
  {four virtual bonds and one physical leg per site};
\node[aqmlabsmall, anchor=north, align=center, text width=4.4cm, text=aqmbroken] at (9.13,-0.34)
  {no 1D order to sweep along, and contracting the network is itself hard: DMRG's local
   eigenproblem has no analogue here};

\end{tikzpicture}

\end{aqmfigure}
